\paragraph{A rigorous time enclosure of the charged trajectory.}
The finite trajectory admits a rigorous error bound on the whole supplied
action-time interval \([0,2]\). The enclosure uses the same charged initial
data and full dressed action as the numerical execution. It controls time
integration on the fixed cone, separately from spatial convergence,
observer provenance and physical clock calibration.

Put \(r=\kappa/|\Omega|=6/(7+3\sqrt5)\), and divide the action by
\(|\Omega|\), which does not change its Euler equations. In the symmetric
sector set \(C=e^{ie\alpha}c\),
\(x=(\operatorname{Re}C,\operatorname{Im}C,
\operatorname{Re}b,\operatorname{Im}b)\), and \(Z(s)=sC+(1-s)b\).
The time kinetic term is
\[
 \tfrac12r\dot\alpha^2+
 \int_0^1 3(1-s)^2\left|\dot Z-ie\dot\alpha sZ\right|^2\,ds.
\]
Thus \(\alpha\) is cyclic. The constant scalar velocity block and
its inverse are
\[
 H=\begin{pmatrix}1/5&3/10\\3/10&6/5\end{pmatrix}\otimes I_2,
 \qquad H^{-1}=\begin{pmatrix}8&-2\\-2&4/3\end{pmatrix}\otimes I_2.
\]
Writing the full kinetic metric as
\(\left(\begin{smallmatrix}d&\ell^{\mathsf T}\\\ell&H\end{smallmatrix}\right)\),
one obtains
\begin{align*}
 \ell&=\frac e{10}\bigl(C_i+b_i,-C_r-b_r,C_i+2b_i,-C_r-2b_r\bigr),\\
 w=H^{-1}\ell&=\bigl(e(3C_i+2b_i)/5,-e(3C_r+2b_r)/5,
                         e(b_i-C_i)/15,e(C_r-b_r)/15\bigr),\\
 D=d-\ell^{\mathsf T}H^{-1}\ell
   &=r+\frac{2e^2}{525}|C-b|^2\ \geq r>\frac25.
\end{align*}
Before volume normalization, this is the positive denominator
\(\kappa+2e^2|\Omega||C-b|^2/525\geq\kappa\).
All coefficients follow from the exact moments
\(\int_0^1 3s^n(1-s)^2ds=6/[(n+1)(n+2)(n+3)]\).

Let \(p\in\mathbb R^4\) be the scalar canonical momenta of the normalized
action. The initial cyclic momentum is exactly zero and remains zero. The
normalized potential is
\[
 W(x)=r|C-b|^2+\tfrac12\langle|Z|^2\rangle
                         +\tfrac18\langle|Z|^4\rangle,
 \qquad \langle f\rangle=\int_0^1 3(1-s)^2f(s)\,ds.
\]
For \(X=|C|^2\), \(Y=\operatorname{Re}(\overline Cb)\), \(B=|b|^2\),
the moments are
\[
 \langle|Z|^2\rangle=X/10+3Y/10+3B/5,
 \qquad
 \langle|Z|^4\rangle=
 \frac{X^2+3XY+2XB+4Y^2+10YB+15B^2}{35}.
\]
The reduced Hamiltonian is
\(\mathcal H=\tfrac12p^{\mathsf T}H^{-1}p+
(w^{\mathsf T}p)^2/(2D)+W\).
Together with cyclic-coordinate reconstruction, its equations are the
rational system
\begin{align}
 a&=-w^{\mathsf T}p/D,&
 \dot\alpha&=a,&\dot x&=H^{-1}p-wa,\nonumber\\
 \dot p_j&=a\,\partial_j(w^{\mathsf T}p)
               +\tfrac12a^2\partial_jD-\partial_jW.
 \label{eq:whitney-charged-rational-ode}
\end{align}
Its nine initial coordinates are
\[
 y(0)=(\alpha,x,p)=(0,1,0,1,0,0,3/10,0,-3/10).
\]
They give \(\dot\alpha(0)=3/(40r)\),
\(\dot C(0)=i(3+e\dot\alpha(0))\) and \(\dot b(0)=-i\),
exactly the preceding charged temporal-gauge initial condition. The
normalized conserved energy is \(49/40+9/(3200r)\).

\begin{proposition}[Uniform time enclosure for the finite charged action]
\label{prop:whitney-charged-time-enclosure}
For \(e=1/4\), \(m^2=1/2\), \(g=1/4\), the exact initial-value solution
of \eqref{eq:whitney-charged-rational-ode} exists uniquely on \([0,2]\).
The supplied piecewise degree-32 polynomial \(P\), on eighty intervals
of exact length \(1/40\), satisfies
\begin{equation}
 \sup_{0\leq t\leq2}\|y(t)-P(t)\|_\infty\leq10^{-20}.
 \label{eq:whitney-charged-time-enclosure}
\end{equation}
At shared interval endpoints either adjacent polynomial may be used;
the approximants need not join exactly. Under the inverse phase map, the
enclosed solution lifts to all 68 real temporal-gauge Euler equations and
all 13 Gauss constraints. At the 81 nominal times \(j/40\), its five
original positions and five original velocities differ from the stored binary64 samples by at most \(10^{-10}\) in each coordinate.
\end{proposition}
\begin{proof}
The Hamiltonian reduction above follows by completing the scalar kinetic
square. Its positive Schur denominator makes the rational vector field
smooth on all real canonical coordinates. The exact initial cyclic momentum
is zero. Canonical equivalence and the full symmetry lift therefore identify
the enclosed solution with the same finite-action solution; the compatible
initial Gauss law propagates by the action's Noether identity.

Here is the numerical certificate argument, including continuous times.
On a step starting from interval enclosure \(Y_j\), the verifier checks a
closed box \(B_j\) and the strict Picard inclusion
\[
 Y_j+[0,h]F(B_j)\subset\operatorname{int}B_j,
 \qquad h=1/40.
\]
Local existence, smoothness and the first-exit argument keep every relevant
solution in this box for the whole step. A contraction estimate is not
required. Formal power-series substitution in \(\dot y=F(y)\) produces
intervals \(A_k(Y_j)\) containing the normalized derivatives
\(y^{(k)}(t_j)/k!\). The same recurrence initialized on the entire box
\(B_j\) bounds \(y^{(33)}(t)/33!\) at every time in the step. Taylor's
theorem consequently gives, for every \(u\in[0,h]\),
\[
 y(t_j+u)\in\sum_{k=0}^{32}A_k(Y_j)u^k+A_{33}(B_j)u^{33}.
\]
At \(u=h\) this interval is checked to lie in the next recorded enclosure.
For each supplied polynomial coefficient \(P_{j,k}\), the verifier also
bounds \(|A_k(Y_j)-P_{j,k}|\) coordinatewise and sums those radii times
\(h^k\), together with the remainder radius times \(h^{33}\).
Every resulting bound is at most \(10^{-20}\). Induction over the eighty
steps proves existence and the displayed uniform error estimate, including
both choices of approximant at a join.

Interval endpoints are integers divided by \(2^{256}\); all elementary
operations round outwards by integer division. The exact rational step is
enclosed rather than rounded to a point. An integer square-root inequality
encloses \(\sqrt5\), and each division checks separation from zero.
Polynomial coefficients are exact multiples of \(2^{-96}\).
The verifier independently reconstructs the scalar mass block, Schur
complement and potential from rational simplex moments, and does not import
the producer or a floating-point ODE solver.

For the historical comparison, each stored binary64 value is interpreted
as its exact dyadic rational. Interval sine and cosine evaluations use
degree-40 Taylor polynomials and the real Lagrange remainder
\(|\theta|^{41}/41!\). Applying these to \(c=e^{-ie\alpha}C\) and its
time derivative encloses the original position and velocity coordinates at
each nominal sample time. Their checked differences give the stated
\(10^{-10}\) bound. The historical samples play no role in constructing or
proving the trajectory enclosure.
\end{proof}

The producer, independent verifier and exact interval data are
\path{code/electromagnetism/whitney_charged_enclosure.py},
\path{code/electromagnetism/verify_whitney_charged_enclosure.py} and
\path{code/electromagnetism/runtime/whitney_charged_enclosure_receipt.json}.
Source identities, coefficient encodings, initial data and interpretation
are part of the checked packet. The enclosure certifies this finite
classical initial-value problem, not the truth of its supplied physical
inputs. It neither calibrates the action parameter in seconds nor certifies
the separate polygonal Jacobi-clock quadrature, charged spatial continuum
limit, observer execution or quantum evolution. The interval argument is
implemented and independently replayed; it is not a Lean formalization.
